\subsubsection{Two structural reductions for six-dimensional quotients}

\begin{lemma}\label{B:six-nonsplit-triangular-row}
Let $A$ be a rational algebra of dimension six and let $V$ be a
two-dimensional module whose action algebra is the full triangular
two-by-two rational matrix algebra.  If $A/\operatorname{rad}A$
is not a product of copies of $\mathbb Q$, then eventual vanishing
of $\operatorname{Tor}^A_n(W,V)$ implies
$\operatorname{Tor}^A_2(W,V)=0$ for every finite-dimensional right
module $W$.
\end{lemma}
\begin{proof}
The nonsplit extension defining $V$ uses two distinct rational
characters.  Thus $A/J$ contains two copies of $\mathbb Q$, and
$J/J^2$ has a nonzero, at least one-dimensional, component joining
their vertices.  Every other simple factor which is not
$\mathbb Q$ has dimension at least two.  A nonzero radical
bimodule incident to a field factor of degree at least two has
dimension at least two; the same lower bound holds for a matrix
factor.  Such a radical component together with the required
rational cross-arrow would exceed the radical dimension available
in a six-dimensional algebra.  A matrix factor of dimension four
already uses all six semisimple dimensions after the two rational
characters and leaves no radical at all.

Consequently every non-rational simple factor is isolated in
$J/J^2$, and therefore in every $J^n/J^{n+1}$ by multiplication.
It splits off as a semisimple direct factor.  The factor acting
nontrivially on $V$ is a nonlocal split basic rational algebra of
dimension at most four.  The nonlocal part of
Theorem~\ref{B:contracted-four-Tor} applies to every such algebra,
independently of a monoid basis: eventual Tor vanishing forces its
degree-two Tor group to vanish.  Direct-factor decomposition gives
the claimed conclusion for $A$.
\end{proof}

\begin{lemma}[An odd Euler characteristic at a local corner]
\label{B:six-local-corner-odd-Euler}
Let $B$ be a four-dimensional local rational algebra with residue
field $k=\mathbb Q$, and put
\[
 A=\begin{pmatrix}k&k\\0&B\end{pmatrix},
\]
where the right $B$-action on the displayed cross bimodule is its
residue character.  Let $V$ be the two-dimensional left module
with one copy of $k$ at each vertex and with the cross-arrow
acting as the identity.  For every five-dimensional right
$A$-module $W$, the groups $\operatorname{Tor}^A_n(W,V)$ are
nonzero in arbitrarily high degrees.
\end{lemma}
\begin{proof}
Write $f,e$ for the first and second vertex idempotents.
A right $A$-module is a triple
\[
 W=(F,U,t),\qquad F=Wf,\quad U=We,\quad
 t:F\otimes_k k\longrightarrow U,
\]
where $U$ is a right $B$-module and $t$ is $B$-linear when its
domain has the residue action.  Put $d=\dim F$, so
$\dim U+d=5$.  There is an explicit exact sequence of right
$A$-modules
\[
 0\longrightarrow(0,k^d)
 \xrightarrow{\ (\iota,-t)\ }
 (F\otimes_k fA)\oplus(0,U)
 \longrightarrow W\longrightarrow0.
\]
The last map is the identity at the $f$-vertex and sends the
$e$-vertex pair $(z,u)$ to $t(z)+u$.  Its kernel is
$z\mapsto(z,-t(z))$, proving exactness and specifying every map.

The module $fA$ is projective and $fA\otimes_A V=fV=k$.
The induction $(0,U)=U\otimes_B eA$ is exact and sends free
$B$-modules to projective $A$-modules, since $eA=(0,B)$.
Consequently
\[
 (0,U)\otimes_A^{\mathbb L}V
 \simeq U\otimes_B^{\mathbb L}k.
\]
Tensor the displayed short exact sequence with $V$ in the derived
category.  Projection away from the summand $F$ induces a map
\[
 W\otimes_A^{\mathbb L}V
 \longrightarrow
 C\otimes_B^{\mathbb L}k,
 \qquad
 C=\operatorname{Cone}(t:k^d\longrightarrow U).
\]
Its fiber is the vector space $F$ in degree zero.  Equivalently,
the source is the cone of
$k^d\otimes_B^{\mathbb L}k\to
 F\oplus(U\otimes_B^{\mathbb L}k)$ with components the
augmentation and $-t$, whereas the target discards the first
component.  Thus this map is a homology isomorphism in every
degree at least two.

Suppose the Tor groups in the statement vanish eventually.
Then $C\otimes_B^{\mathbb L}k$ has homology in only finitely
many degrees.  A bounded complex of finite-dimensional modules
over a finite-dimensional local algebra has a minimal free
resolution, constructed successively by projective covers.
Its differential entries lie in the radical.  Tensoring with
$k$ therefore makes every differential zero; the assumed
bounded homology forces this minimal resolution to be bounded.
Hence $C$ is a perfect complex of $B$-modules.

Every finite projective $B$-module is free and has rational
dimension divisible by four.  The rational Euler characteristic
of a perfect $B$-complex is therefore divisible by four.  But
the explicit two-term complex $C$ has Euler characteristic
\[
 \chi_{\mathbb Q}(C)=\dim U-d=5-2d,
\]
which is odd.  This contradiction proves the assertion.
\end{proof}

\begin{remark}\label{B:six-local-corner-both-orientations}
For the opposite triangular orientation
$\left(\begin{smallmatrix}B&k\\0&k\end{smallmatrix}\right)$,
the two-dimensional nonsplit cross-arrow module is itself the
projective at the $k$-vertex.  Its higher Tor groups are zero.
Thus the two orientations of a local four-dimensional corner
with a single radical-annihilated cross-arrow cannot produce
Tor dimensions $2,0,1,0,\ldots$ with augmentation dimensions
$(\dim W,\dim V)=(5,2)$.
\end{remark}
